\uuid{ZpOj}
\titre{Équations complexes}
\auteur{bodin ; hanani ; mzari}
\organisation{exo7}
\licence{CC BY-NC-SA 4.0}
\datecreate{2020-01-20}
\niveau{L1}
\module{Algèbre}
\chapitre{Nombres complexes}
\sousChapitre{Racine carrée, équation du second degré}
\difficulte{1}
\theme{Équations, équations complexes}
\exo7qcmid{lille-1.nombres-complexes.018}

\contenu{
\begin{question}
\qclassification{104.02,104.03,104.04}
\qtags{motcle=équations complexes}
Les racines cubiques de $1+i$ sont :
\begin{answers}
    \bad{$z_k=\sqrt[3]{2}e^{i(\frac{\pi}{12}+\frac{2k\pi}{3})}, k=0,1,2$}
    \good{$z_k=\sqrt[6]{2}e^{i(\frac{\pi}{12}+\frac{2k\pi}{3})}, k=0,1,2$}

    \good{$z_k=\sqrt[6]{2}e^{i(\frac{\pi}{12}-\frac{2k\pi}{3})}, k=0,1,2$ }

    \bad{$z_k=\sqrt[3]{2}e^{i(\frac{\pi}{12}-\frac{2k\pi}{3})}, k=0,1,2$}
\end{answers}
\begin{explanations}
On résoud l'équation  : $z^3=1+i= \sqrt 2e^{i\frac{\pi}{4}}$.

\end{explanations}
\end{question}
}
